% GENERATED FILE. Edit canonical lesson frontmatter, then run npm run generate:books.
% canonical-entries: 146

\chapter*{Glossary}
\addcontentsline{toc}{chapter}{Glossary}
\markboth{Glossary}{Glossary}

This glossary collects every word and phrase taught in the book. Pronunciation appears when it differs from the written form; chapter numbers show where each entry is introduced.

\noindent\begin{minipage}[t]{\linewidth}
\raggedright
\textbf{abrir}\par
\small to open — a regular -ir verb with one irregular part, the participle abierto\par
\footnotesize Introduced in Chapter 149.
\end{minipage}\par\medskip

\noindent\begin{minipage}[t]{\linewidth}
\raggedright
\textbf{adiós}\par
\small goodbye\par
\footnotesize Introduced in Chapter 12.
\end{minipage}\par\medskip

\noindent\begin{minipage}[t]{\linewidth}
\raggedright
\textbf{amarillo}\par
\small yellow — from \textquotedblleft{}a little bitter,\textquotedblright{} not from any Latin word for yellow at all\par
\footnotesize Introduced in Chapter 134.
\end{minipage}\par\medskip

\noindent\begin{minipage}[t]{\linewidth}
\raggedright
\textbf{así que}\par
\small thus, in that way — and with que after it, \textquotedblleft{}so\textquotedblright{}; built on the same sīc that gave you sí\par
\footnotesize Introduced in Chapter 158.
\end{minipage}\par\medskip

\noindent\begin{minipage}[t]{\linewidth}
\raggedright
\textbf{ayudar}\par
\small to help — literally the same word as English \textquotedblleft{}aid\textquotedblright{}, arriving by a different road\par
\footnotesize Introduced in Chapter 144.
\end{minipage}\par\medskip

\noindent\begin{minipage}[t]{\linewidth}
\raggedright
\textbf{azul}\par
\small blue — an Arabic loan that became the plain everyday word, where French's cousin stayed poetic\par
\footnotesize Introduced in Chapter 122.
\end{minipage}\par\medskip

\noindent\begin{minipage}[t]{\linewidth}
\raggedright
\textbf{beber}\par
\small to drink (another -er verb — cementing the pattern)\par
\footnotesize Introduced in Chapter 32.
\end{minipage}\par\medskip

\noindent\begin{minipage}[t]{\linewidth}
\raggedright
\textbf{bien}\par
\small well / good / fine\par
\footnotesize Introduced in Chapter 1.
\end{minipage}\par\medskip

\noindent\begin{minipage}[t]{\linewidth}
\raggedright
\textbf{buenas noches}\par
\small good evening / good night\par
\footnotesize Introduced in Chapter 2.
\end{minipage}\par\medskip

\noindent\begin{minipage}[t]{\linewidth}
\raggedright
\textbf{buenas tardes}\par
\small good afternoon\par
\footnotesize Introduced in Chapter 2.
\end{minipage}\par\medskip

\noindent\begin{minipage}[t]{\linewidth}
\raggedright
\textbf{buenos días}\par
\small good morning (literally \textquotedblleft{}good days\textquotedblright{})\par
\footnotesize Introduced in Chapter 1.
\end{minipage}\par\medskip

\noindent\begin{minipage}[t]{\linewidth}
\raggedright
\textbf{café}\par
\small coffee / a café — one useful noun before the polite request that follows\par
\footnotesize Introduced in Chapter 14.
\end{minipage}\par\medskip

\noindent\begin{minipage}[t]{\linewidth}
\raggedright
\textbf{caminar}\par
\small to walk — built on a Celtic word for a path, beside andar, whose own past is unsettled\par
\footnotesize Introduced in Chapter 148.
\end{minipage}\par\medskip

\noindent\begin{minipage}[t]{\linewidth}
\raggedright
\textbf{cerrar}\par
\small to close — from the bar dropped across a door, not from the Latin verb English inherited\par
\footnotesize Introduced in Chapter 149.
\end{minipage}\par\medskip

\noindent\begin{minipage}[t]{\linewidth}
\raggedright
\textbf{comer}\par
\small to eat — your first -er verb, and the family whose yo-form you already own\par
\footnotesize Introduced in Chapter 30.
\end{minipage}\par\medskip

\noindent\begin{minipage}[t]{\linewidth}
\raggedright
\textbf{comí}\par
\small the regular singular -er preterite — comí, comiste, comió\par
\footnotesize Introduced in Chapter 87.
\end{minipage}\par\medskip

\noindent\begin{minipage}[t]{\linewidth}
\raggedright
\textbf{cómo}\par
\small how — a whole Latin phrase, quo modo, worn smooth into one word\par
\footnotesize Introduced in Chapter 6.
\end{minipage}\par\medskip

\noindent\begin{minipage}[t]{\linewidth}
\raggedright
\textbf{comprar}\par
\small to buy — from Latin's \textquotedblleft{}get ready, procure\textquotedblright{}, and not from the verb behind English compare\par
\footnotesize Introduced in Chapter 157.
\end{minipage}\par\medskip

\noindent\begin{minipage}[t]{\linewidth}
\raggedright
\textbf{conocer}\par
\small to meet, to be acquainted with — a verb built out of the moment knowing begins\par
\footnotesize Introduced in Chapter 154.
\end{minipage}\par\medskip

\noindent\begin{minipage}[t]{\linewidth}
\raggedright
\textbf{conseguir}\par
\small to get, to obtain — getting as following a thing until you have it\par
\footnotesize Introduced in Chapter 152.
\end{minipage}\par\medskip

\noindent\begin{minipage}[t]{\linewidth}
\raggedright
\textbf{contar}\par
\small to tell a story — and also to count, because in Spanish, as once in English, a story is something you reckon up\par
\footnotesize Introduced in Chapter 150.
\end{minipage}\par\medskip

\noindent\begin{minipage}[t]{\linewidth}
\raggedright
\textbf{contestar}\par
\small to answer — a Roman courtroom word that English turned into a fight and Spanish turned into a reply\par
\footnotesize Introduced in Chapter 157.
\end{minipage}\par\medskip

\noindent\begin{minipage}[t]{\linewidth}
\raggedright
\textbf{correr}\par
\small to run — a plain -er verb sitting on one of the largest roots English ever borrowed\par
\footnotesize Introduced in Chapter 148.
\end{minipage}\par\medskip

\noindent\begin{minipage}[t]{\linewidth}
\raggedright
\textbf{creer}\par
\small to believe, to think — putting your heart somewhere, where pensar weighs it up\par
\footnotesize Introduced in Chapter 158.
\end{minipage}\par\medskip

\noindent\begin{minipage}[t]{\linewidth}
\raggedright
\textbf{cuando}\par
\small \textquotedblleft{}when\textquotedblright{} — the word that puts a scene and an event in one sentence, so the two pasts can be seen doing different jobs\par
\footnotesize Introduced in Chapter 95.
\end{minipage}\par\medskip

\noindent\begin{minipage}[t]{\linewidth}
\raggedright
\textbf{de nada}\par
\small you're welcome (literally \textquotedblleft{}of nothing\textquotedblright{})\par
\footnotesize Introduced in Chapter 7.
\end{minipage}\par\medskip

\noindent\begin{minipage}[t]{\linewidth}
\raggedright
\textbf{deber}\par
\small should, ought to — and to owe, because the word is literally \textquotedblleft{}to have something from someone\textquotedblright{}\par
\footnotesize Introduced in Chapter 158.
\end{minipage}\par\medskip

\noindent\begin{minipage}[t]{\linewidth}
\raggedright
\textbf{decir}\par
\small to say or tell — with the singular present forms digo, dices, and dice\par
\footnotesize Introduced in Chapter 51.
\end{minipage}\par\medskip

\noindent\begin{minipage}[t]{\linewidth}
\raggedright
\textbf{día}\par
\small day (el día — masculine)\par
\footnotesize Introduced in Chapter 1.
\end{minipage}\par\medskip

\noindent\begin{minipage}[t]{\linewidth}
\raggedright
\textbf{dieciséis — diecinueve}\par
\small 16-19 — not words to memorise but a formula, diez y X, with the seam still visible\par
\footnotesize Introduced in Chapter 132.
\end{minipage}\par\medskip

\noindent\begin{minipage}[t]{\linewidth}
\raggedright
\textbf{dónde}\par
\small where?\par
\footnotesize Introduced in Chapter 32.
\end{minipage}\par\medskip

\noindent\begin{minipage}[t]{\linewidth}
\raggedright
\textbf{dormir}\par
\small to sleep — and the other half of the boot, where a stressed o breaks to ue\par
\footnotesize Introduced in Chapter 148.
\end{minipage}\par\medskip

\noindent\begin{minipage}[t]{\linewidth}
\raggedright
\textbf{e $\to$ ie and o $\to$ ue}\par
\small comparing the two vowel changes inside the learned singular forms\par
\footnotesize Introduced in Chapter 50.
\end{minipage}\par\medskip

\noindent\begin{minipage}[t]{\linewidth}
\raggedright
\textbf{el / la}\par
\small the (masculine / feminine)\par
\footnotesize Introduced in Chapter 1.
\end{minipage}\par\medskip

\noindent\begin{minipage}[t]{\linewidth}
\raggedright
\textbf{el agua}\par
\small water — a word that barely changed from Latin, and takes el without ever stopping being feminine\par
\footnotesize Introduced in Chapter 126.
\end{minipage}\par\medskip

\noindent\begin{minipage}[t]{\linewidth}
\raggedright
\textbf{el hermano, la hermana}\par
\small brother and sister — NOT from frater/soror, but from an adjective meaning \textquotedblleft{}of the same seed\textquotedblright{}\par
\footnotesize Introduced in Chapter 123.
\end{minipage}\par\medskip

\noindent\begin{minipage}[t]{\linewidth}
\raggedright
\textbf{el libro}\par
\small book — named after the part of a tree people wrote on before paper\par
\footnotesize Introduced in Chapter 59.
\end{minipage}\par\medskip

\noindent\begin{minipage}[t]{\linewidth}
\raggedright
\textbf{el padre, la madre}\par
\small father and mother — among the oldest reconstructible words in the whole family\par
\footnotesize Introduced in Chapter 123.
\end{minipage}\par\medskip

\noindent\begin{minipage}[t]{\linewidth}
\raggedright
\textbf{el pan}\par
\small bread — and the compañero who shares it with you\par
\footnotesize Introduced in Chapter 126.
\end{minipage}\par\medskip

\noindent\begin{minipage}[t]{\linewidth}
\raggedright
\textbf{el tiempo}\par
\small the one noun that covers both \textquotedblleft{}time\textquotedblright{} and \textquotedblleft{}weather\textquotedblright{} — because Latin tempus already did\par
\footnotesize Introduced in Chapter 131.
\end{minipage}\par\medskip

\noindent\begin{minipage}[t]{\linewidth}
\raggedright
\textbf{el vino}\par
\small wine — one Latin root that reached English by three separate roads, plus one English built at home\par
\footnotesize Introduced in Chapter 126.
\end{minipage}\par\medskip

\noindent\begin{minipage}[t]{\linewidth}
\raggedright
\textbf{entender}\par
\small to understand — literally \textquotedblleft{}to stretch toward\textquotedblright{}; the second e$\to$ie stem-changer\par
\footnotesize Introduced in Chapter 137.
\end{minipage}\par\medskip

\noindent\begin{minipage}[t]{\linewidth}
\raggedright
\textbf{escribir}\par
\small to write — from a Latin verb meaning \textquotedblleft{}to scratch\textquotedblright{}, and the reason Spanish puts an e- in front\par
\footnotesize Introduced in Chapter 139.
\end{minipage}\par\medskip

\noindent\begin{minipage}[t]{\linewidth}
\raggedright
\textbf{español}\par
\small \textquotedblleft{}Spanish\textquotedblright{} — the language, the people, and the Roman name for the peninsula underneath it\par
\footnotesize Introduced in Chapter 22.
\end{minipage}\par\medskip

\noindent\begin{minipage}[t]{\linewidth}
\raggedright
\textbf{esperar}\par
\small to wait — and to hope, and to expect, all in one verb\par
\footnotesize Introduced in Chapter 157.
\end{minipage}\par\medskip

\noindent\begin{minipage}[t]{\linewidth}
\raggedright
\textbf{está en…}\par
\small is in or at… — locate someone with estar and the already-known en\par
\footnotesize Introduced in Chapter 44.
\end{minipage}\par\medskip

\noindent\begin{minipage}[t]{\linewidth}
\raggedright
\textbf{estar}\par
\small to be — but the \emph{standing} kind, straight from Latin stāre\par
\footnotesize Introduced in Chapter 8.
\end{minipage}\par\medskip

\noindent\begin{minipage}[t]{\linewidth}
\raggedright
\textbf{estás / está}\par
\small what estar is for — a current state or location — and the only two forms you need this chapter\par
\footnotesize Introduced in Chapter 8.
\end{minipage}\par\medskip

\noindent\begin{minipage}[t]{\linewidth}
\raggedright
\textbf{estudiar}\par
\small to study (a third -ar verb — the pattern is automatic now)\par
\footnotesize Introduced in Chapter 21.
\end{minipage}\par\medskip

\noindent\begin{minipage}[t]{\linewidth}
\raggedright
\textbf{explicar}\par
\small to explain — literally to unfold, and the verb that names what this whole chapter has been building\par
\footnotesize Introduced in Chapter 158.
\end{minipage}\par\medskip

\noindent\begin{minipage}[t]{\linewidth}
\raggedright
\textbf{fui}\par
\small the shared singular preterite of ser and ir — fui, fuiste, fue\par
\footnotesize Introduced in Chapter 86.
\end{minipage}\par\medskip

\noindent\begin{minipage}[t]{\linewidth}
\raggedright
\textbf{gato}\par
\small cat — a word that travelled out of Africa with the animal, pushing Latin's own fēlēs aside\par
\footnotesize Introduced in Chapter 133.
\end{minipage}\par\medskip

\noindent\begin{minipage}[t]{\linewidth}
\raggedright
\textbf{gracias}\par
\small thank you\par
\footnotesize Introduced in Chapter 7.
\end{minipage}\par\medskip

\noindent\begin{minipage}[t]{\linewidth}
\raggedright
\textbf{gustar}\par
\small to please — the verb English calls \textquotedblleft{}to like\textquotedblright{}, built the other way round\par
\footnotesize Introduced in Chapter 146.
\end{minipage}\par\medskip

\noindent\begin{minipage}[t]{\linewidth}
\raggedright
\textbf{hablar}\par
\small to speak / to talk — at root, \textquotedblleft{}to tell stories\textquotedblright{}\par
\footnotesize Introduced in Chapter 16.
\end{minipage}\par\medskip

\noindent\begin{minipage}[t]{\linewidth}
\raggedright
\textbf{Hablo español}\par
\small your first sentence built from parts — and the rule that languages take no article\par
\footnotesize Introduced in Chapter 22.
\end{minipage}\par\medskip

\noindent\begin{minipage}[t]{\linewidth}
\raggedright
\textbf{hace calor, hace frío, hace sol}\par
\small the three weather words Spanish reports with hacer — \textquotedblleft{}it makes heat,\textquotedblright{} \textquotedblleft{}it makes cold,\textquotedblright{} \textquotedblleft{}it makes sun\textquotedblright{}\par
\footnotesize Introduced in Chapter 131.
\end{minipage}\par\medskip

\noindent\begin{minipage}[t]{\linewidth}
\raggedright
\textbf{hacer}\par
\small to do or make — with the singular present forms hago, haces, and hace\par
\footnotesize Introduced in Chapter 51.
\end{minipage}\par\medskip

\noindent\begin{minipage}[t]{\linewidth}
\raggedright
\textbf{hasta}\par
\small until / up to\par
\footnotesize Introduced in Chapter 12.
\end{minipage}\par\medskip

\noindent\begin{minipage}[t]{\linewidth}
\raggedright
\textbf{hasta luego}\par
\small see you later (literally \textquotedblleft{}until later\textquotedblright{})\par
\footnotesize Introduced in Chapter 13.
\end{minipage}\par\medskip

\noindent\begin{minipage}[t]{\linewidth}
\raggedright
\textbf{hasta mañana}\par
\small see you tomorrow (literally \textquotedblleft{}until tomorrow\textquotedblright{})\par
\footnotesize Introduced in Chapter 13.
\end{minipage}\par\medskip

\noindent\begin{minipage}[t]{\linewidth}
\raggedright
\textbf{hasta pronto}\par
\small see you soon (literally \textquotedblleft{}until soon\textquotedblright{})\par
\footnotesize Introduced in Chapter 13.
\end{minipage}\par\medskip

\noindent\begin{minipage}[t]{\linewidth}
\raggedright
\textbf{hay}\par
\small \textquotedblleft{}there is\textquotedblright{} and \textquotedblleft{}there are\textquotedblright{} — one word for both, which no other Spanish verb form manages\par
\footnotesize Introduced in Chapter 174.
\end{minipage}\par\medskip

\noindent\begin{minipage}[t]{\linewidth}
\raggedright
\textbf{hola}\par
\small hello / hi\par
\footnotesize Introduced in Chapter 1.
\end{minipage}\par\medskip

\noindent\begin{minipage}[t]{\linewidth}
\raggedright
\textbf{ir}\par
\small to go — with the three singular present forms voy, vas, and va\par
\footnotesize Introduced in Chapter 45.
\end{minipage}\par\medskip

\noindent\begin{minipage}[t]{\linewidth}
\raggedright
\textbf{jugar}\par
\small to play — from Latin's word for joking, and the only verb in the language that breaks u to ue\par
\footnotesize Introduced in Chapter 153.
\end{minipage}\par\medskip

\noindent\begin{minipage}[t]{\linewidth}
\raggedright
\textbf{la}\par
\small \textquotedblleft{}it\textquotedblright{} for something feminine — the one word in this pair that never had to change\par
\footnotesize Introduced in Chapter 67.
\end{minipage}\par\medskip

\noindent\begin{minipage}[t]{\linewidth}
\raggedright
\textbf{la cabeza}\par
\small the head — Spanish KEPT Latin's real word, unlike French, which replaced it with slang for \textquotedblleft{}pot\textquotedblright{}\par
\footnotesize Introduced in Chapter 124.
\end{minipage}\par\medskip

\noindent\begin{minipage}[t]{\linewidth}
\raggedright
\textbf{la casa}\par
\small house, home — and a word French wore down so far it stopped looking like a noun\par
\footnotesize Introduced in Chapter 58.
\end{minipage}\par\medskip

\noindent\begin{minipage}[t]{\linewidth}
\raggedright
\textbf{la comida}\par
\small food, a meal — a noun you can build yourself out of a verb you already own\par
\footnotesize Introduced in Chapter 60.
\end{minipage}\par\medskip

\noindent\begin{minipage}[t]{\linewidth}
\raggedright
\textbf{la hora}\par
\small hour / o'clock — telling the time, and a word Spanish inherited almost unchanged from Latin\par
\footnotesize Introduced in Chapter 130.
\end{minipage}\par\medskip

\noindent\begin{minipage}[t]{\linewidth}
\raggedright
\textbf{la mano}\par
\small the hand — feminine despite ending in -o, Spanish's most famous gender exception\par
\footnotesize Introduced in Chapter 124.
\end{minipage}\par\medskip

\noindent\begin{minipage}[t]{\linewidth}
\raggedright
\textbf{la noche}\par
\small night / evening (la noche — feminine)\par
\footnotesize Introduced in Chapter 2.
\end{minipage}\par\medskip

\noindent\begin{minipage}[t]{\linewidth}
\raggedright
\textbf{la primavera, el verano, el otoño, el invierno}\par
\small the four seasons — three straightforward, and one whose MEANING quietly shifted from \textquotedblleft{}spring\textquotedblright{} to \textquotedblleft{}summer\textquotedblright{}\par
\footnotesize Introduced in Chapter 125.
\end{minipage}\par\medskip

\noindent\begin{minipage}[t]{\linewidth}
\raggedright
\textbf{la tarde}\par
\small afternoon (la tarde — feminine; also: late)\par
\footnotesize Introduced in Chapter 2.
\end{minipage}\par\medskip

\noindent\begin{minipage}[t]{\linewidth}
\raggedright
\textbf{le}\par
\small \textquotedblleft{}to him, to her, to you\textquotedblright{} — a Latin case Spanish kept and English threw away\par
\footnotesize Introduced in Chapter 76.
\end{minipage}\par\medskip

\noindent\begin{minipage}[t]{\linewidth}
\raggedright
\textbf{leer}\par
\small to read — from a Latin verb meaning \textquotedblleft{}to gather, to pick out\textquotedblright{}\par
\footnotesize Introduced in Chapter 138.
\end{minipage}\par\medskip

\noindent\begin{minipage}[t]{\linewidth}
\raggedright
\textbf{levantarse}\par
\small to get up, to stand up — \textquotedblleft{}to raise oneself\textquotedblright{}, built from a participle exactly as sentarse was\par
\footnotesize Introduced in Chapter 149.
\end{minipage}\par\medskip

\noindent\begin{minipage}[t]{\linewidth}
\raggedright
\textbf{llamo}\par
\small I call (from llamar, to call)\par
\footnotesize Introduced in Chapter 4.
\end{minipage}\par\medskip

\noindent\begin{minipage}[t]{\linewidth}
\raggedright
\textbf{llovía / llovió}\par
\small the same rain in two tenses — how a story uses the imperfect for the scene and the preterite for the events\par
\footnotesize Introduced in Chapter 150.
\end{minipage}\par\medskip

\noindent\begin{minipage}[t]{\linewidth}
\raggedright
\textbf{llueve}\par
\small rain refuses the hacer pattern and brings its own verb — and its ll- is the sound change you already met in llamar\par
\footnotesize Introduced in Chapter 131.
\end{minipage}\par\medskip

\noindent\begin{minipage}[t]{\linewidth}
\raggedright
\textbf{lo}\par
\small \textquotedblleft{}it\textquotedblright{} — a word you have been saying since your fourth chapter, in a job you have not met yet\par
\footnotesize Introduced in Chapter 66.
\end{minipage}\par\medskip

\noindent\begin{minipage}[t]{\linewidth}
\raggedright
\textbf{lo siento}\par
\small I'm sorry — literally \textquotedblleft{}I feel it,\textquotedblright{} a naming of the feeling rather than a claim about yourself\par
\footnotesize Introduced in Chapter 120.
\end{minipage}\par\medskip

\noindent\begin{minipage}[t]{\linewidth}
\raggedright
\textbf{los meses}\par
\small the months — Roman gods and emperors, marching almost unchanged from Latin\par
\footnotesize Introduced in Chapter 128.
\end{minipage}\par\medskip

\noindent\begin{minipage}[t]{\linewidth}
\raggedright
\textbf{luego}\par
\small then, next — a word you already own from hasta luego, now doing a second job: joining two events in order\par
\footnotesize Introduced in Chapter 150.
\end{minipage}\par\medskip

\noindent\begin{minipage}[t]{\linewidth}
\raggedright
\textbf{lunes, martes, miércoles, jueves, viernes}\par
\small the weekdays (Monday–Friday) — worn-down planet-god names\par
\footnotesize Introduced in Chapter 121.
\end{minipage}\par\medskip

\noindent\begin{minipage}[t]{\linewidth}
\raggedright
\textbf{más despacio}\par
\small \textquotedblleft{}more slowly\textquotedblright{} — the request that makes the last sentence solvable\par
\footnotesize Introduced in Chapter 15.
\end{minipage}\par\medskip

\noindent\begin{minipage}[t]{\linewidth}
\raggedright
\textbf{me}\par
\small me / myself\par
\footnotesize Introduced in Chapter 4.
\end{minipage}\par\medskip

\noindent\begin{minipage}[t]{\linewidth}
\raggedright
\textbf{me llamo}\par
\small my name is (literally \textquotedblleft{}I call myself\textquotedblright{})\par
\footnotesize Introduced in Chapter 4.
\end{minipage}\par\medskip

\noindent\begin{minipage}[t]{\linewidth}
\raggedright
\textbf{mediodía, medianoche}\par
\small noon and midnight — both halves still fully alive as ordinary Spanish words, unlike French's worn-down midi/minuit\par
\footnotesize Introduced in Chapter 129.
\end{minipage}\par\medskip

\noindent\begin{minipage}[t]{\linewidth}
\raggedright
\textbf{mi, tu, su}\par
\small my, your informal, and your formal or his or her — before a singular noun\par
\footnotesize Introduced in Chapter 45.
\end{minipage}\par\medskip

\noindent\begin{minipage}[t]{\linewidth}
\raggedright
\textbf{mucho}\par
\small much / a lot\par
\footnotesize Introduced in Chapter 6.
\end{minipage}\par\medskip

\noindent\begin{minipage}[t]{\linewidth}
\raggedright
\textbf{mucho gusto}\par
\small pleased to meet you (literally \textquotedblleft{}much pleasure\textquotedblright{})\par
\footnotesize Introduced in Chapter 6.
\end{minipage}\par\medskip

\noindent\begin{minipage}[t]{\linewidth}
\raggedright
\textbf{negro, blanco}\par
\small black and white — the two words, and the straight-line Latin history of the first one\par
\footnotesize Introduced in Chapter 122.
\end{minipage}\par\medskip

\noindent\begin{minipage}[t]{\linewidth}
\raggedright
\textbf{no}\par
\small no / not\par
\footnotesize Introduced in Chapter 7.
\end{minipage}\par\medskip

\noindent\begin{minipage}[t]{\linewidth}
\raggedright
\textbf{no entiendo}\par
\small \textquotedblleft{}I don't understand\textquotedblright{} — the most useful sentence a beginner owns\par
\footnotesize Introduced in Chapter 15.
\end{minipage}\par\medskip

\noindent\begin{minipage}[t]{\linewidth}
\raggedright
\textbf{nos}\par
\small \textquotedblleft{}us\textquotedblright{} — and the word hiding inside the verb ending you have used since hablamos\par
\footnotesize Introduced in Chapter 71.
\end{minipage}\par\medskip

\noindent\begin{minipage}[t]{\linewidth}
\raggedright
\textbf{nuestro, nuestra}\par
\small our — matching a known masculine or feminine singular noun\par
\footnotesize Introduced in Chapter 50.
\end{minipage}\par\medskip

\noindent\begin{minipage}[t]{\linewidth}
\raggedright
\textbf{oír}\par
\small to hear — the half of listening you do not choose, and a verb that sprouts a y\par
\footnotesize Introduced in Chapter 148.
\end{minipage}\par\medskip

\noindent\begin{minipage}[t]{\linewidth}
\raggedright
\textbf{once — quince}\par
\small 11-15 — five single opaque words, each one a Latin compound worn down until the seam disappeared\par
\footnotesize Introduced in Chapter 132.
\end{minipage}\par\medskip

\noindent\begin{minipage}[t]{\linewidth}
\raggedright
\textbf{os}\par
\small \textquotedblleft{}you all\textquotedblright{} as the object — the shortest word in the set, and the one that stops at a border\par
\footnotesize Introduced in Chapter 73.
\end{minipage}\par\medskip

\noindent\begin{minipage}[t]{\linewidth}
\raggedright
\textbf{para}\par
\small \textquotedblleft{}for\textquotedblright{} pointing forward — at a destination, a purpose, a deadline, a person meant to receive something\par
\footnotesize Introduced in Chapter 159.
\end{minipage}\par\medskip

\noindent\begin{minipage}[t]{\linewidth}
\raggedright
\textbf{pensar}\par
\small to think — at root \textquotedblleft{}to weigh\textquotedblright{}, and a stem-changer that breaks e$\to$ie\par
\footnotesize Introduced in Chapter 136.
\end{minipage}\par\medskip

\noindent\begin{minipage}[t]{\linewidth}
\raggedright
\textbf{perdón}\par
\small the other sorry — \textquotedblleft{}give completely\textquotedblright{} — for fault and for getting past people, where lo siento is for regret\par
\footnotesize Introduced in Chapter 120.
\end{minipage}\par\medskip

\noindent\begin{minipage}[t]{\linewidth}
\raggedright
\textbf{perro}\par
\small dog — a genuinely unsolved etymology, and the Latin word it pushed out\par
\footnotesize Introduced in Chapter 133.
\end{minipage}\par\medskip

\noindent\begin{minipage}[t]{\linewidth}
\raggedright
\textbf{poder}\par
\small to be able or can — with the singular present forms puedo, puedes, and puede\par
\footnotesize Introduced in Chapter 50.
\end{minipage}\par\medskip

\noindent\begin{minipage}[t]{\linewidth}
\raggedright
\textbf{poner}\par
\small to put or place — with the singular present forms pongo, pones, and pone\par
\footnotesize Introduced in Chapter 52.
\end{minipage}\par\medskip

\noindent\begin{minipage}[t]{\linewidth}
\raggedright
\textbf{por}\par
\small \textquotedblleft{}for\textquotedblright{} pointing back or through — at a cause, a route, an exchange; and you have been saying it since your first page\par
\footnotesize Introduced in Chapter 160.
\end{minipage}\par\medskip

\noindent\begin{minipage}[t]{\linewidth}
\raggedright
\textbf{por favor}\par
\small please (literally \textquotedblleft{}for {[}a{]} favour\textquotedblright{})\par
\footnotesize Introduced in Chapter 14.
\end{minipage}\par\medskip

\noindent\begin{minipage}[t]{\linewidth}
\raggedright
\textbf{porque}\par
\small because — two tiny words fused into one, and the reason a sequence becomes a story\par
\footnotesize Introduced in Chapter 150.
\end{minipage}\par\medskip

\noindent\begin{minipage}[t]{\linewidth}
\raggedright
\textbf{preguntar}\par
\small to ask (a question) — and the reason English's one \textquotedblleft{}ask\textquotedblright{} becomes two Spanish verbs\par
\footnotesize Introduced in Chapter 143.
\end{minipage}\par\medskip

\noindent\begin{minipage}[t]{\linewidth}
\raggedright
\textbf{qué}\par
\small what? (a new question word for verbs you already know)\par
\footnotesize Introduced in Chapter 32.
\end{minipage}\par\medskip

\noindent\begin{minipage}[t]{\linewidth}
\raggedright
\textbf{querer}\par
\small to want — with the singular present forms quiero, quieres, and quiere\par
\footnotesize Introduced in Chapter 50.
\end{minipage}\par\medskip

\noindent\begin{minipage}[t]{\linewidth}
\raggedright
\textbf{regular}\par
\small so-so / not great (also \textquotedblleft{}más o menos\textquotedblright{})\par
\footnotesize Introduced in Chapter 9.
\end{minipage}\par\medskip

\noindent\begin{minipage}[t]{\linewidth}
\raggedright
\textbf{rojo}\par
\small red — a cousin of French rouge and English red, but by a different branch of the same very old family\par
\footnotesize Introduced in Chapter 122.
\end{minipage}\par\medskip

\noindent\begin{minipage}[t]{\linewidth}
\raggedright
\textbf{sábado, domingo}\par
\small the weekend (Saturday, Sunday) — where religion overwrote astronomy\par
\footnotesize Introduced in Chapter 121.
\end{minipage}\par\medskip

\noindent\begin{minipage}[t]{\linewidth}
\raggedright
\textbf{salir}\par
\small to leave or go out — with the singular present forms salgo, sales, and sale\par
\footnotesize Introduced in Chapter 53.
\end{minipage}\par\medskip

\noindent\begin{minipage}[t]{\linewidth}
\raggedright
\textbf{se llama}\par
\small he/she calls himself; you (formal) call yourself\par
\footnotesize Introduced in Chapter 6.
\end{minipage}\par\medskip

\noindent\begin{minipage}[t]{\linewidth}
\raggedright
\textbf{seis, siete, ocho, nueve, diez}\par
\small the numbers 6–10 and the old counting hidden in month names\par
\footnotesize Introduced in Chapter 43.
\end{minipage}\par\medskip

\noindent\begin{minipage}[t]{\linewidth}
\raggedright
\textbf{sentarse}\par
\small to sit down — literally \textquotedblleft{}to seat oneself\textquotedblright{}, the reflexive that lands the action back on you\par
\footnotesize Introduced in Chapter 149.
\end{minipage}\par\medskip

\noindent\begin{minipage}[t]{\linewidth}
\raggedright
\textbf{ser}\par
\small to be — the verb used here to identify someone or say what someone is\par
\footnotesize Introduced in Chapter 44.
\end{minipage}\par\medskip

\noindent\begin{minipage}[t]{\linewidth}
\raggedright
\textbf{ser vs estar}\par
\small identify with ser; describe a current state or location with estar\par
\footnotesize Introduced in Chapter 44.
\end{minipage}\par\medskip

\noindent\begin{minipage}[t]{\linewidth}
\raggedright
\textbf{sí}\par
\small yes\par
\footnotesize Introduced in Chapter 7.
\end{minipage}\par\medskip

\noindent\begin{minipage}[t]{\linewidth}
\raggedright
\textbf{soy de…}\par
\small I am from… — use ser with de to state origin\par
\footnotesize Introduced in Chapter 44.
\end{minipage}\par\medskip

\noindent\begin{minipage}[t]{\linewidth}
\raggedright
\textbf{te quiero}\par
\small \textquotedblleft{}you\textquotedblright{} as the object — and the sentence it makes, which is the one people actually need\par
\footnotesize Introduced in Chapter 68.
\end{minipage}\par\medskip

\noindent\begin{minipage}[t]{\linewidth}
\raggedright
\textbf{tener}\par
\small to have — a common verb with two visible changes\par
\footnotesize Introduced in Chapter 43.
\end{minipage}\par\medskip

\noindent\begin{minipage}[t]{\linewidth}
\raggedright
\textbf{tengo · hago · digo}\par
\small comparing the three learned yo forms that end in -go\par
\footnotesize Introduced in Chapter 51.
\end{minipage}\par\medskip

\noindent\begin{minipage}[t]{\linewidth}
\raggedright
\textbf{tomar}\par
\small to take (and, across Latin America, to drink) — an everyday word with no settled origin\par
\footnotesize Introduced in Chapter 142.
\end{minipage}\par\medskip

\noindent\begin{minipage}[t]{\linewidth}
\raggedright
\textbf{trabajar}\par
\small to work (a second -ar verb — same pattern)\par
\footnotesize Introduced in Chapter 20.
\end{minipage}\par\medskip

\noindent\begin{minipage}[t]{\linewidth}
\raggedright
\textbf{traer}\par
\small to bring — Latin's word for dragging, and a verb with a direction built into it\par
\footnotesize Introduced in Chapter 151.
\end{minipage}\par\medskip

\noindent\begin{minipage}[t]{\linewidth}
\raggedright
\textbf{tú / usted}\par
\small two of the three words Spanish uses to one person — and the one English quietly retired\par
\footnotesize Introduced in Chapter 5.
\end{minipage}\par\medskip

\noindent\begin{minipage}[t]{\linewidth}
\raggedright
\textbf{tú o usted}\par
\small the choice English lost — and the surprise that usted takes he/she verb forms\par
\footnotesize Introduced in Chapter 5.
\end{minipage}\par\medskip

\noindent\begin{minipage}[t]{\linewidth}
\raggedright
\textbf{uno, dos, tres, cuatro, cinco}\par
\small the numbers 1–5\par
\footnotesize Introduced in Chapter 43.
\end{minipage}\par\medskip

\noindent\begin{minipage}[t]{\linewidth}
\raggedright
\textbf{veinte}\par
\small twenty — the ceiling of the teens, and a word that was worn but never fused\par
\footnotesize Introduced in Chapter 132.
\end{minipage}\par\medskip

\noindent\begin{minipage}[t]{\linewidth}
\raggedright
\textbf{venir}\par
\small to come — with the singular present forms vengo, vienes, and viene\par
\footnotesize Introduced in Chapter 54.
\end{minipage}\par\medskip

\noindent\begin{minipage}[t]{\linewidth}
\raggedright
\textbf{ver}\par
\small to see — veo, ves, ve\par
\footnotesize Introduced in Chapter 91.
\end{minipage}\par\medskip

\noindent\begin{minipage}[t]{\linewidth}
\raggedright
\textbf{verde}\par
\small green — from a verb meaning \textquotedblleft{}to flourish,\textquotedblright{} and a cousin of verdant but not of English green\par
\footnotesize Introduced in Chapter 134.
\end{minipage}\par\medskip

\noindent\begin{minipage}[t]{\linewidth}
\raggedright
\textbf{viví}\par
\small the regular singular -ir preterite — viví, viviste, vivió\par
\footnotesize Introduced in Chapter 87.
\end{minipage}\par\medskip

\noindent\begin{minipage}[t]{\linewidth}
\raggedright
\textbf{vivir}\par
\small to live — your first -ir verb, and a family that asks you to learn nothing new\par
\footnotesize Introduced in Chapter 31.
\end{minipage}\par\medskip

\noindent\begin{minipage}[t]{\linewidth}
\raggedright
\textbf{vos}\par
\small \textquotedblleft{}you\textquotedblright{} across a large part of the Spanish-speaking world — to recognise, not to say\par
\footnotesize Introduced in Chapter 5.
\end{minipage}\par\medskip

\noindent\begin{minipage}[t]{\linewidth}
\raggedright
\textbf{voy a + infinitive}\par
\small I am going to… — the near future with ir, a, and a known infinitive\par
\footnotesize Introduced in Chapter 45.
\end{minipage}\par\medskip

\noindent\begin{minipage}[t]{\linewidth}
\raggedright
\textbf{y / ¿y tú?}\par
\small and / and you? — a tiny connector that hands the conversation back\par
\footnotesize Introduced in Chapter 9.
\end{minipage}\par\medskip

\noindent\begin{minipage}[t]{\linewidth}
\raggedright
\textbf{¿cómo está usted?}\par
\small how are you? — assemble the formal question from cómo, estar, and usted\par
\footnotesize Introduced in Chapter 8.
\end{minipage}\par\medskip

\noindent\begin{minipage}[t]{\linewidth}
\raggedright
\textbf{¿cómo se dice?}\par
\small the single most useful sentence a learner owns — and it is this chapter's grammar, doing exactly what the chapter said\par
\footnotesize Introduced in Chapter 166.
\end{minipage}\par\medskip

\noindent\begin{minipage}[t]{\linewidth}
\raggedright
\textbf{¿cómo se llama usted?}\par
\small what's your name? (formal)\par
\footnotesize Introduced in Chapter 6.
\end{minipage}\par\medskip

\noindent\begin{minipage}[t]{\linewidth}
\raggedright
\textbf{¿Cuántos años tienes?}\par
\small how old are you? — literally, how many years do you have?\par
\footnotesize Introduced in Chapter 43.
\end{minipage}\par\medskip
